\documentclass[a4paper,11pt]{article}

\usepackage[margin=2.5cm]{geometry}
\usepackage{amsmath,amssymb}
\usepackage{graphicx}
\usepackage{hyperref}
\usepackage{cite}
\usepackage{algorithm}
\usepackage{algorithmicx}
\usepackage{algpseudocode}
\usepackage{booktabs}
\usepackage{enumitem}

\title{Spotless: Grid-Free Deconvolution\\for Radio Interferometry}
\author{Tim Molteno\\
\texttt{tim@elec.ac.nz}}
\date{\today}

\begin{document}

\maketitle

\begin{abstract}
Spotless is a point-source deconvolution algorithm for radio interferometry
that operates directly in visibility space, without gridding. By exploiting
Parseval's theorem, the total energy in the image can be computed from the
visibilities without performing a Fourier transform. Spotless builds a model
of the sky as a sum of point sources by minimising the residual energy. This
approach handles non-coplanar baselines and all-sky fields of view without
requiring $w$-projection or other wide-field corrections. Two variants are
presented: sequential spotless (analogous to CLEAN) and MultiSpotless, which
jointly optimises all model components.
\end{abstract}

%=============================================================================
\section{Introduction}
%=============================================================================

Radio interferometers measure the complex visibility function
$V(u, v, w)$, which is related to the sky brightness distribution
$I(l, m)$ through the van~Cittert--Zernike theorem~\cite{thompson2017}:
%
\begin{equation}
V(u, v, w) = \iint \frac{I(l, m)}{\sqrt{1 - l^2 - m^2}}
             \, e^{-2\pi j [ul + vm + w(\sqrt{1-l^2-m^2} - 1)]}
             \, dl \, dm.
\label{eq:vcz}
\end{equation}
%
Here $(l, m, n)$ are direction cosines relative to the phase centre, and
$(u, v, w)$ are baseline coordinates in wavelengths.

Traditional imaging pipelines grid the visibilities onto a regular
$(u,v)$-plane and apply a Fast Fourier Transform (FFT) to produce a
dirty image. The CLEAN algorithm~\cite{hogbom1974} then deconvolves the
dirty image by iteratively subtracting scaled copies of the dirty beam.
Variants such as the Cotton--Schwab algorithm~\cite{schwab1984} perform
the subtraction in visibility space to reduce aliasing artefacts.

However, gridding introduces several complications: $w$-terms cause
distortions for non-coplanar arrays; all-sky imaging requires careful
treatment of the horizon; and wide-field imaging suffers from
direction-dependent effects that are difficult to model on a
regular grid.

Spotless avoids gridding entirely by working directly with the
visibilities. The key insight is Parseval's theorem: the total energy
(power) of an image can be computed from the visibilities without
performing a Fourier transform. This enables a model-fitting approach
where point sources are located and characterised by minimising the
residual visibility energy.

%=============================================================================
\section{Theory}
%=============================================================================

\subsection{Point-Source Visibility}

A point source at direction cosines $(l, m)$ with brightness $A$
produces a visibility:
%
\begin{equation}
V_P(u, v, w) = A^2 \, e^{-2\pi j [ul + vm + w(n-1)]},
\label{eq:ptvis}
\end{equation}
%
where $n = \sqrt{1 - l^2 - m^2}$. In the Spotless implementation,
the source position is parameterised in elevation--azimuth coordinates
$(\theta, \phi)$ rather than $(l, m)$, with the conversion:
%
\begin{align}
l &= \sin\phi \cos\theta, \\
m &= \cos\phi \cos\theta, \\
n &= \sin\theta.
\label{eq:lmn}
\end{align}

%-----------------------------------------------------------------------------
\subsection{Energy in Visibility Space}
%-----------------------------------------------------------------------------

For discrete visibilities $v_k$ measured at $M$ baseline points,
Parseval's theorem for the Discrete Fourier Transform states:
%
\begin{equation}
\sum_{n=0}^{N-1} |x_n|^2 = \frac{1}{N} \sum_{k=0}^{N-1} |X_k|^2,
\end{equation}
%
where $x_n$ is the image and $X_k$ is its Fourier transform. Since the
visibilities are samples of the Fourier transform of the sky brightness,
the total energy $E(V)$ in the image is proportional to the sum of
squared visibility amplitudes:
%
\begin{equation}
E(V) \propto \sum_{k=0}^{M-1} v_k v_k^*.
\label{eq:energy}
\end{equation}
%
No gridding or inverse Fourier transform is required to compute this
quantity---it is evaluated directly from the visibility data.

%-----------------------------------------------------------------------------
\subsection{Sequential Spotless}
%-----------------------------------------------------------------------------

The sequential Spotless algorithm builds a model of $N$ point sources by
iteratively finding and subtracting the brightest remaining source.
Beginning with the measured visibilities $V^{(0)} = V$, each step
$i = 0, 1, \dots$ solves:
%
\begin{equation}
(A_i, \theta_i, \phi_i) =
\underset{A,\theta,\phi}{\arg\min}\;
E\bigl(V^{(i)} - A^2 V_P(\theta, \phi)\bigr),
\label{eq:step}
\end{equation}
%
where $V_P(\theta, \phi)$ is the unit-amplitude point-source visibility
from equation~\eqref{eq:ptvis}. The source is added to the model and
subtracted from the residual:
%
\begin{align}
\mathcal{M} &\leftarrow \mathcal{M} \cup \{(A_i, \theta_i, \phi_i)\}, \\
V^{(i+1)} &= V^{(i)} - A_i^2 V_P(\theta_i, \phi_i).
\end{align}
%
The algorithm terminates when the residual energy stops decreasing,
i.e., when $E(V^{(i+1)}) \geq E(V^{(i)})$.

The cost function in equation~\eqref{eq:step} is a smooth function of
three parameters $(A, \theta, \phi)$. The optimisation is performed using
the L-BFGS-B algorithm~\cite{byrd1995}, a quasi-Newton method that
supports bound constraints and uses a limited-memory Hessian
approximation. An initial guess is obtained by imaging the current
residual onto a HEALPix sphere~\cite{gorski2005} and locating the
brightest pixel.

%-----------------------------------------------------------------------------
\subsection{MultiSpotless}
%-----------------------------------------------------------------------------

MultiSpotless extends the sequential approach by jointly re-optimising
all model components at each step. After $i$ sources have been found,
the $(i+1)$-th step solves the full $3(i+1)$-dimensional problem:
%
\begin{equation}
\underset{\{A_k,\theta_k,\phi_k\}_{k=0}^{i}}
{\arg\min}\;
E\bigl(V - \sum_{k=0}^{i} A_k^2 V_P(\theta_k, \phi_k)\bigr).
\label{eq:multistep}
\end{equation}
%
This allows existing sources to be refined as new ones are added,
potentially yielding a more accurate model at the cost of increased
computational complexity. Because the full parameter space mixes
amplitudes and angles at very different scales, MultiSpotless employs
the Nelder--Mead simplex algorithm~\cite{nelder1965} with adaptive
parameters, which is derivative-free and robust to poor scaling.

%=============================================================================
\section{Algorithm}
%=============================================================================

\begin{algorithm}[htbp]
\caption{Sequential Spotless}
\label{alg:spotless}
\begin{algorithmic}[1]
\State \textbf{Input:} Visibilities $V$, field-of-view sphere $\mathcal{S}$
\State $\mathcal{M} \leftarrow \emptyset$
\State $V_{\text{res}} \leftarrow V$
\State $E_0 \leftarrow \infty$
\Repeat
    \State Image $V_{\text{res}}$ onto $\mathcal{S}$
    \State $(A_0, \theta_0, \phi_0) \leftarrow \text{peak of image}$
    \State $x_0 \leftarrow [A_0, \theta_0, \phi_0]$
    \State $x^* \leftarrow \arg\min_x E(V_{\text{res}} - V_P(x))$
    \State $(A, \theta, \phi) \leftarrow x^*$
    \State $E_1 \leftarrow E(V_{\text{res}} - A^2 V_P(\theta, \phi))$
    \If{$E_1 \geq E_0$} \textbf{break} \EndIf
    \State $\mathcal{M} \leftarrow \mathcal{M} \cup \{(A, \theta, \phi)\}$
    \State $V_{\text{res}} \leftarrow V_{\text{res}} - A^2 V_P(\theta, \phi)$
    \State $E_0 \leftarrow E_1$
\Until{max iterations}
\State \textbf{return} $\mathcal{M}$, $V_{\text{res}}$
\end{algorithmic}
\end{algorithm}

The MultiSpotless variant replaces the single-source optimisation
(line~8) with a joint optimisation over all model parameters as in
equation~\eqref{eq:multistep}.

%=============================================================================
\section{Implementation}
%=============================================================================

Spotless is implemented in Python and distributed as the
\texttt{spotless} package. Key design features include:

\begin{itemize}[leftmargin=*,nosep]
    \item \textbf{Grid-free operation.} All computations are performed
          directly on the visibility data; no gridding or degridding steps
          are required.
    \item \textbf{All-sky support.} The field of view is represented as a
          HEALPix sphere, enabling natural handling of wide-field and
          all-sky observations.
    \item \textbf{Non-coplanar baselines.} The full $(u, v, w)$ baseline
          coordinates are used in the point-source visibility calculation
          (equation~\ref{eq:ptvis}), avoiding the need for $w$-projection.
    \item \textbf{Plug-and-play data sources.} Spotless accepts
          visibilities from a TART telescope API, from JSON snapshot
          files, or from CASA measurement sets (via \texttt{disko}).
    \item \textbf{Multiple output formats.} Reconstructed images can be
          exported as FITS, PNG, SVG, PDF, or HDF5.
\end{itemize}

The initial peak search and image reconstruction use the HEALPix
pixelisation scheme. Beam smoothing during reconstruction is performed
with the \texttt{healpy} spherical harmonic transform.

%=============================================================================
\section{Results}
%=============================================================================

Figure~\ref{fig:results} shows a comparison between a dirty image and
a Spotless reconstruction from TART telescope data. Spotless produces a
clean point-source model with substantially reduced sidelobe artefacts.

\begin{figure}[htbp]
\centering
\includegraphics[width=0.45\textwidth]{../img/gridless.png}
\hfill
\includegraphics[width=0.45\textwidth]{../img/spotless.png}
\caption{Left: dirty (gridless) image. Right: Spotless reconstruction
         showing the deconvolved point-source model.}
\label{fig:results}
\end{figure}

%=============================================================================
\section{Discussion}
%=============================================================================

\subsection{Advantages over CLEAN}

Spotless offers several advantages over traditional CLEAN-based methods:
%
\begin{enumerate}[leftmargin=*,nosep]
    \item No gridding artefacts---the model is fit directly to the
          visibility data.
    \item No requirement for a regular $(u,v)$ sampling.
    \item Natural handling of non-coplanar baselines via the full
          $(u, v, w)$ phase term.
    \item All-sky fields of view are handled without special treatment.
    \item The model is a compact parametric representation (list of point
          sources) rather than a pixel map.
\end{enumerate}

\subsection{Limitations and Future Work}

\begin{itemize}[leftmargin=*,nosep]
    \item \textbf{Point-source assumption.} Spotless models the sky as a
          collection of point sources. Extended emission is not well
          represented. Future work could incorporate Gaussian or
          shapelet source models.
    \item \textbf{Antenna beam.} The current implementation assumes an
          isotropic hemispherical antenna response. Incorporating a
          direction-dependent beam model would improve accuracy,
          particularly at low elevations.
    \item \textbf{Uncertainty quantification.} Running a Markov Chain
          Monte Carlo (MCMC) sampler on the MultiSpotless model would
          provide uncertainty estimates for source parameters and a
          natural stopping criterion based on posterior precision.
    \item \textbf{Computational cost.} MultiSpotless scales as
          $\mathcal{O}(N^2)$ in the number of sources due to the joint
          optimisation. Block-coordinate descent or alternating
          optimisation could reduce this cost.
\end{itemize}

%=============================================================================
\section{Conclusion}
%=============================================================================

Spotless provides a grid-free approach to radio interferometric
deconvolution by exploiting Parseval's theorem to compute image energy
directly from visibilities. By fitting a point-source model through
energy minimisation, it avoids the gridding and degridding steps
required by traditional imaging pipelines. The sequential variant offers
speed comparable to CLEAN, while MultiSpotless provides improved
accuracy through joint optimisation. The method is particularly
well-suited to all-sky imaging with non-coplanar arrays.

%=============================================================================
\bibliographystyle{plain}
\bibliography{spotless}

\end{document}
